% latexmk -cd -lualatex works/A_19/appendix.tex
% latexmk -c -cd works/A_19/appendix.tex

\documentclass{ees}

\newlist{lyricslist}{description}{1}
\setlist[lyricslist]{
  partopsep=0pt,
  labelindent=0em,
  labelwidth=6em,
  leftmargin=6em,
  labelsep=0pt,
  first=\ltseries,
  font=\normalfont\itshape\ltseries,
  style=multiline
}

\newenvironment{lyrics}[1]{%
  \subsection{#1}\nopagebreak%
  \begin{lyricslist}%
  \let\voice\item%
}{%
  \end{lyricslist}%
}

\begin{document}

\clearscrheadfoot

\chapter{Lyrics}

\textit{Unterredende Perſonen.}\\
Jahel, die beherzte Heldin. Canto.\\
Debora, Eine mit Prophetiſchen Geiſt begaabte Richterin in Israel. Alto. \\
Sisera, der Feldfürſt des Cananiter Königs Jabin. Tenore.\\
Barack, der Israelitiſchen Feld-Obriſter. Baſso.\\
Chorus. der ſtreitenden Israelitern.

\section{Actus primus}

\begin{lyrics}{Scena prima}
  \voice[Chorus der Israeliten]
  Gott Israel, blickh auf daß Joch,\\
  ſo unßern Ohnmachtsnaken drüket.\\
  Erbarme dich und helffe doch\\
  eh unßer Hofnungsgeiſt erſtiket.\\
  Zernicht den Zwang von Bänden\\
  der bangen Dienſtbarkeit\\
  und mach von Feundesbanden\\
  dein Israel befreyt,\\
  und mach dein Volk befreyt.
\end{lyrics}

\begin{lyrics}{Scena secunda}
  \voice[Debora]
  Er ſterbe, er ſterbe. Diſer Schluß,\\
  der unßre Freyheit flüglen muß,\\
  ſoll unſrer Feunde Dämpfer ſeyn.

  \voice[Barak]
  Welch Raßen nimmt dich ein?\\
  Laßtu den Sanfftmuethsgeiſt\\
  auß ſeinen Trieben gleitten,\\
  du wüetteſt mit dir ſelbſt,\\
  waß will mir daß bedeuthen?\\
  Wohin verlierſt du dich?

  \voice[Debora]
  Ja Barak höre mich,\\
  und durch mich Gottes Willen,\\
  den Ich und du erfüllen\\
  durch ſein Geheißen ſoll.\\
  Der Herr, von dem daß Wohl\\
  ſowie daß Weh entſprünget,\\
  der Helden höbt und ſtürzt,\\
  der Wueth und Lüſt bezwünget,\\
  der unſern Feund erhub\\
  damit er tieffer fallt,\\
  will nun durch unſer Schwerd,\\
  gelenket durch den Gewald,\\
  der ßeine Allmacht wird\\
  in unſre Kräfften ſenken,\\
  durch den geſtürzten Feund\\
  unß unſer Freyheit ſchenkhen.

  \voice[Barak]
  Wie dunkel redeſt du\\
  von Troſt ergänzten Sachen,\\
  waß ſoll ich mir hierauß\\
  dan vor Begriffe machen.

  \voice[Debora]
  Held, höre fehrner zue,\\
  du wirſt ſogleich daß Liecht\\
  in meinen Reden fünden,\\
  den Werkzeug, welchen Gott\\
  zur Geißel unſrer Sünden\\
  bißhero angewant,\\
  ſoll unſre Tapferkeit\\
  alß zweyter Werkzeug auch\\
  ſelbſt wider brechen,\\
  und zwar in einen Streit,\\
  den du mit Sisera ſolſt wagen,\\
  wo du zwar Israel wirſt rächen,\\
  doch ſeinen Feund nicht ſchlagen.

  \voice[Barak]
  Du ſprichſt geheimnußvoll\\
  und lokerſt meinen Sinn\\
  zu ganz verwirten Denken,\\
  wohin wird endlich ſich\\
  dein klarer Ausſpruch lenken?

  \voice[Debora]
  Ich rede, wie ich ſoll,\\
  nun frage weiter nicht,\\
  die Zeit wird dich ſchon lehren,\\
  waß Sieg und Schlagen heiſt,\\
  jez folge nun dem Herren,\\
  der dich durch beſondern Ruff\\
  zu ſeines Israels\\
  Befreyer hat ernennet,\\
  du ſolſt die Kriegesflamm,\\
  ſo um den Tabor brennet,\\
  mit tapferen Behuf\\
  von zehentauſend Mann\\
  der Kinder Sabulon\\
  und Nephtalim erſtüken.\\
  Doch wird die Siegeskron\\
  von dißesmahl ſich nicht\\
  um deine Schläffe ſchicken,\\
  dan Gott wird Sisera\\
  und ſein erboßtes Leben\\
  in eines Weibes Hände geben.

  \voice[Barak]
  Bringt eines Weibes Hand\\
  ſo villes wohl zu Stand,\\
  daß ſie auch Helden überwindet?

  \voice[Debora]
  Wenn Gott die Siegeskrafft\\
  auf Weiber Armbe bindet,\\
  ſo zwüngt ein Weibe diß,\\
  waß auch bey Helden offt\\
  ſchon unerzwünglich ſcheint und hieß.
\end{lyrics}

\begin{lyrics}{Aria prima}
  \voice[Debora]
  Es muß nicht ſtäts ein Donnerſtein\\
  den ſtolzen Keder ſchmettern,\\
  und nicht alzeit ein Sturmwind ſeyn\\
  die Äeſte abzublättern.\\
  Indeme auch zuweil\\
  ein ſchwach doch ſcharffes Beill\\
  ihn khan zu Boden werffen.\\[1ex]
  Alßo wird ein ſchwaches Weib,\\
  welchen Gott die Hände führet,\\
  wan ſie auch den ſtärkeſten Leib\\
  eines Helden nur berühret\\
  den Sieg erwarthen därffen.
\end{lyrics}

% \begin{lyrics}{}
%   \voice[]
% \end{lyrics}

% \begin{lyrics}{}
%   \voice[]
% \end{lyrics}

% \begin{lyrics}{}
%   \voice[]
% \end{lyrics}

% \begin{lyrics}{}
%   \voice[]
% \end{lyrics}

% \begin{lyrics}{}
%   \voice[]
% \end{lyrics}

% \begin{lyrics}{}
%   \voice[]
% \end{lyrics}

% \begin{lyrics}{}
%   \voice[]
% \end{lyrics}

% \begin{lyrics}{}
%   \voice[]
% \end{lyrics}

% \begin{lyrics}{}
%   \voice[]
% \end{lyrics}

% \begin{lyrics}{}
%   \voice[]
% \end{lyrics}

% \begin{lyrics}{}
%   \voice[]
% \end{lyrics}

% \begin{lyrics}{}
%   \voice[]
% \end{lyrics}

% \begin{lyrics}{}
%   \voice[]
% \end{lyrics}

% \begin{lyrics}{}
%   \voice[]
% \end{lyrics}

% \begin{lyrics}{}
%   \voice[]
% \end{lyrics}

% \begin{lyrics}{}
%   \voice[]
% \end{lyrics}

% \begin{lyrics}{}
%   \voice[]
% \end{lyrics}

% \begin{lyrics}{}
%   \voice[]
% \end{lyrics}

% \begin{lyrics}{}
%   \voice[]
% \end{lyrics}

% \begin{lyrics}{}
%   \voice[]
% \end{lyrics}

\end{document}
